\scanpage{159}
\mspage{104}
il suffit de prouver que les deux flèches sont égales \uline{dans} $(\underline{\mathrm{Fib}}\,S)W_S^{-1}$. Cet énoncé est indépendant de la donnée de $\mathcal L$, et on a envie de le prouver.

Peut-être y a-t-il moyen d'\unsure{expliciter|exploiter} un langage géométrique, pour s'y reconnaître. \deleted{Pour \unsure{un objet}} Identifions les objets \inserted{$X$} de $\mathcal M$ aux foncteurs $\mathcal M^\circ\longrightarrow(\mathrm{Ens})$ qu'ils définissent,
\[
X\longmapsto\bigl(T\longmapsto X(T)=\mathrm{Hom}_{\mathcal M}(T,X)\bigr)\text{.}
\]
Alors, pour tt $T$
\[
\mathrm{Ch}(X)(T)
\]
est un ensemble, appelé l'\uline{\deleted{\unsure{ens}}ensemble des chemins dans} $X(T)$. Il s'envoie (par $p_X$) fonctoriellement \inserted{(en $T$)} dans $(X\times X)(T)=X(T)\times X(T)$, i.e. chaque chemin \inserted{$\unsure{c|\gamma}$} dans $X(T)$ \unsure{à|a} une origine \inserted{$\sigma(c)$}, et un but $\beta(c)$. De plus, par $\alpha_X$, $X(T)$ s'\unsure{envoit|envoie} (fonctoriellement en $T$) sur l'ens. des chemins de $X(T)$ : pour chaque élément \inserted{$x$} de $X(T)$, il y a donc un chemin (noté \deleted{$\unsure{\mathrm{id}}$}$\unsure{1}_x$ ou $\alpha_x$) \deleted{\unsure{de}} bien défini de source et but $x$, qu'on appelle le chemin \unsure{trivial|trivide}, ou vide. Ceci dit,
% archive sheet number: 158
